\documentclass[11pt]{article}
\usepackage[a4paper,margin=1in]{geometry}
\usepackage[T1]{fontenc}
\usepackage[utf8]{inputenc}
\usepackage{lmodern}
\usepackage{amsmath,amssymb,amsthm,mathtools}
\usepackage{booktabs}
\usepackage{microtype}
\usepackage[hidelinks]{hyperref}
\usepackage{enumitem}
\usepackage{array}

\newtheorem{theorem}{Theorem}
\newtheorem{lemma}{Lemma}
\newtheorem{proposition}{Proposition}
\newtheorem{corollary}{Corollary}
\theoremstyle{definition}
\newtheorem{definition}{Definition}
\theoremstyle{remark}
\newtheorem{remark}{Remark}

\newcommand{\Acal}{\mathcal A}
\newcommand{\tauac}{\tau_{\rm ac}}

\title{Exact Tightness on the Second Subset Diagonal\\of the Avoid-or-Compress Lemma}
\author{Ryutaro Yonezu\\Independent Researcher}
\date{22 September 2026\\Revised draft -- not peer reviewed}

\begin{document}
\maketitle

\begin{abstract}
Let $\Acal=(Q,\Sigma,\delta)$ be a synchronizing deterministic finite automaton, let $\varnothing\ne A\subsetneq S\subseteq Q$, and call a word $w$ successful if either $A\nsubseteq S w$ or $|S w|<|S|$. The avoid-or-compress lemma gives a successful word of length at most $n-|A|$, where $n=|Q|$. We determine the exact worst-case value on the second subset diagonal $|S|-|A|=2$. If $S=Q$, the exact value is $1$. If $S\subsetneq Q$, the bound $n-|A|=n-|S|+2$ is sharp, even after simultaneously requiring a binary alphabet, strong connectivity, and synchronization. We give explicit extremal families for every proper-subset parameter pair and prove exactness by a safe-subset invariant. Strong connectivity is immediate from the construction, while synchronization follows from a cycle-and-adjacent-merge argument. An independent breadth-first verifier checks 1,140 parameter pairs with $1\le |A|\le30$ and $3\le n-|A|\le40$, with no failures. Thus the general avoid-or-compress upper bound is exactly tight along the entire proper second diagonal under the binary strongly connected synchronizing restrictions.
\end{abstract}

\section{Introduction}
Avoiding and compressing words are two of the basic local mechanisms used in the theory of synchronizing automata. Szyku\l a's avoid-or-compress lemma states that, whenever a proper target subset $A\subsetneq S$ can eventually be avoided from the image of $S$, one can either avoid $A$ or compress $S$ within $n-|A|$ steps \cite{Szykula2018,Szykula2026}. The recent open-problems survey emphasizes that tight parameter-dependent bounds are largely unknown, especially when one insists that extremal witnesses remain synchronizing and strongly connected \cite{Szykula2026}.

This paper closes the first nontrivial proper diagonal beyond the fixed-state case. We impose
\[
|S|-|A|=2
\]
and determine the exact worst-case length. The answer is particularly clean: except at the full-set boundary $S=Q$, the general upper bound $n-|A|$ is attained exactly by an explicit binary, strongly connected, synchronizing family.

The constructions are simple enough that the lower bound can be proved by tracking only the two states of $S\setminus A$. The parity of $L=n-|A|$ determines whether the extremal word is $b^{L-1}a$ or $b^{L-2}a^2$. The synchronization proof is independent of the avoid-or-compress lower bound: after a short entry word the full state set lies in a core on which one letter acts as a full cycle, while a second word induces one adjacent merge. Conjugating that merge around the cycle yields a reset word.

\section{Preliminaries}
We write the action of a word on the right. Thus $q w$ is the state reached from $q$ after reading $w$, and for $X\subseteq Q$,
\[
Xw=\{qw:q\in X\}.
\]

\begin{definition}[Avoid-or-compress distance]
Let $\varnothing\ne A\subsetneq S\subseteq Q$. Define
\[
\tauac(\Acal;S,A)=\min\{|w|: A\nsubseteq Sw\text{ or }|Sw|<|S|\}.
\]
\end{definition}

A word counted by $\tauac$ is called \emph{successful}; all shorter words are \emph{safe}. The avoid-or-compress lemma gives the upper bound
\begin{equation}\label{eq:aoc-upper}
\tauac(\Acal;S,A)\le n-|A|
\end{equation}
whenever the hypotheses of the lemma hold; in particular this applies in the synchronizing setting considered here \cite{Szykula2018,Szykula2026}.

For the second diagonal write
\[
r=|A|,\qquad |S|=r+2,\qquad L=n-r=n-|A|.
\]
For a proper subset $S\subsetneq Q$ we have $L\ge3$.

\section{Main theorem}
\begin{theorem}[Exact second diagonal]\label{thm:main}
Fix integers $n$ and $s$ with $3\le s\le n$, and let $|A|=s-2$.
\begin{enumerate}[label=(\roman*)]
\item If $s=n$, then every synchronizing $n$-state DFA satisfies
\[
\tauac(\Acal;Q,A)=1,
\]
and this value is exact.
\item If $s<n$, then every synchronizing $n$-state DFA satisfies
\[
\tauac(\Acal;S,A)\le n-s+2,
\]
and for every such parameter pair there exists a binary, strongly connected, synchronizing automaton with
\[
\tauac(\Acal;S,A)=n-s+2.
\]
\end{enumerate}
Consequently the general upper bound $n-|A|$ is exactly tight on the entire proper second diagonal, even under the binary, strongly connected, and synchronizing restrictions.
\end{theorem}

The upper bound is \eqref{eq:aoc-upper}. The rest of the paper constructs and verifies the matching lower bound.

\section{Extremal construction}
Fix $r\ge1$ and $L\ge3$. Let
\[
A=\{A_0,\ldots,A_{r-1}\},\qquad
B=\{B_0,\ldots,B_{L-1}\},
\]
\[
Q=A\cup B,\qquad S=A\cup\{B_0,B_1\}.
\]
The alphabet is $\{a,b\}$. Letter $b$ fixes $A$ and cyclically permutes $B$:
\[
A_i b=A_i,\qquad B_j b=B_{j+1\bmod L}.
\]
On $A$, letter $a$ is the chain
\[
A_0\to A_1\to\cdots\to A_{r-1}\to B_0.
\]
Its action on $B$ depends on the parity of $L$.

If $L$ is odd, set
\begin{equation}\label{eq:odd-a}
B_j a=\begin{cases}
A_0,&j\text{ even},\\
B_1,&j\text{ odd}.
\end{cases}
\end{equation}
If $L$ is even, set
\begin{equation}\label{eq:even-a}
B_j a=\begin{cases}
A_0,&j\text{ even},\\
B_1,&j\text{ odd and }j<L-1,\\
B_2,&j=L-1.
\end{cases}
\end{equation}

\section{Exact avoid-or-compress distance}
\begin{lemma}[Odd $L$ safe corridor]\label{lem:odd}
Assume $L$ is odd. Every word of length less than $L$ is safe, while
\[
w=b^{L-1}a
\]
is successful. Hence $\tauac(\Acal;S,A)=L$.
\end{lemma}
\begin{proof}
For $0\le k\le L-1$ define
\[
S_k=A\cup\{B_k,B_{k+1\bmod L}\}.
\]
We have $S b^k=S_k$. If $0\le k\le L-2$, then $k$ and $k+1$ have opposite parity. By \eqref{eq:odd-a}, the two $B$-states are sent to $\{A_0,B_1\}$, while the $A$-chain contributes
\[
\{A_1,\ldots,A_{r-1},B_0\}.
\]
Thus $S_k a=S$ for $k\le L-2$.

Every word can be written as alternating blocks of $b$'s separated by occurrences of $a$. Consider any occurrence of $a$ reached from $S$ after a block $b^k$ with $0\le k\le L-2$. The identity $S_k a=S$ shows that the whole block $b^k a$ returns the active subset exactly to $S$. Hence, in a shortest successful word, every such regular-corridor block may be deleted without changing the subset from which the remaining suffix starts. Repeating this deletion removes every occurrence of $a$ before the wrap-around defect. Thus any shortest successful word must begin with at least $L-1$ consecutive copies of $b$ before its first effective $a$, so no word of length less than $L$ can succeed. After exactly $L-1$ applications of $b$, the active pair is $\{B_{L-1},B_0\}$. Since $L$ is odd, both indices are even, so $a$ maps both states to $A_0$. The image loses one state and is compressed. Thus $b^{L-1}a$ is successful of length $L$.
\end{proof}

\begin{lemma}[Even $L$ safe corridor]\label{lem:even}
Assume $L$ is even. Every word of length less than $L$ is safe, while
\[
w=b^{L-2}a^2
\]
is successful. Hence $\tauac(\Acal;S,A)=L$.
\end{lemma}
\begin{proof}
For $0\le k\le L-3$, the consecutive pair $\{B_k,B_{k+1}\}$ contains one even state and one ordinary odd state. Equation \eqref{eq:even-a} therefore gives $S_k a=S$, exactly as in the odd case. Thus every occurrence of $a$ before the final two-state defect simply returns the active subset to $S$.

After $L-2$ applications of $b$, the active pair is $\{B_{L-2},B_{L-1}\}$. The first $a$ maps it to $\{A_0,B_2\}$, so
\[
S b^{L-2}a=A\cup\{B_0,B_2\}.
\]
This image still contains all of $A$ and still has size $r+2$, hence is safe. A second $a$ sends both $B_0$ and $B_2$ to $A_0$, causing compression. Therefore $b^{L-2}a^2$ is successful and has length $L$.

As in the odd case, write a word as alternating blocks of $b$'s separated by occurrences of $a$. Whenever an $a$ is applied after $b^k$ with $0\le k\le L-3$, the identity $S_k a=S$ shows that the block $b^k a$ returns exactly to $S$. Such a regular-corridor block can therefore be deleted from any shortest successful word. After all such deletions, a shortest candidate must first reach the defect directly from $S$, which requires $L-2$ consecutive copies of $b$. The next letter $a$ occurs at total length $L-1$ and gives $A\cup\{B_0,B_2\}$, which is still safe. Consequently every successful word needs at least one further letter and hence has length at least $L$. Since $b^{L-2}a^2$ succeeds at length $L$, no shorter successful word exists.
\end{proof}

Lemmas \ref{lem:odd} and \ref{lem:even} prove the lower-bound part of Theorem \ref{thm:main} for every proper subset.

\section{Strong connectivity and synchronization}
\begin{proposition}[Strong connectivity]\label{prop:sc}
Every automaton in the construction is strongly connected.
\end{proposition}
\begin{proof}
From every $A_i$, repeated applications of $a$ reach $B_0$. Letter $b$ then reaches every $B_j$. From any even-indexed $B_j$, letter $a$ reaches $A_0$. Hence every state reaches $A_0$, and $A_0$ reaches every state.
\end{proof}

For synchronization define the core
\[
C=A\cup\{B_0\},\qquad m=|C|=r+1.
\]
On $C$, letter $a$ is the $m$-cycle
\[
A_0\to A_1\to\cdots\to A_{r-1}\to B_0\to A_0.
\]
Write $\alpha=a|_C$.

\begin{lemma}[Entry into the core]\label{lem:entry}
If $L$ is odd, the word
\[
w_0=ab^{L-1}a
\]
maps $Q$ into $C$. If $L$ is even, the word
\[
w_0=a^2b^{L-2}a^2
\]
maps $Q$ into $C$.
\end{lemma}
\begin{proof}
This follows by direct substitution from \eqref{eq:odd-a} and \eqref{eq:even-a}. The first one or two $a$'s collapse all $B$-states into a constant-size set, the power of $b$ positions that set at the unique wrap defect, and the final $a$ or $a^2$ places every image in $A\cup\{B_0\}$.
\end{proof}

\begin{lemma}[Adjacent merge on the core]\label{lem:merge}
There is a word $\gamma$ whose restriction to $C$ fixes every core state except one state, which is mapped to its predecessor on the $\alpha$-cycle.
\end{lemma}
\begin{proof}
For odd $L$ set
\[
u=bab^{L-1}a,
\]
and for even $L$ set
\[
u=b^{L-1}a^2.
\]
A direct evaluation on $C$ shows that $u$ acts as $\alpha^2$ except that the final two positions of the $\alpha$-cycle have the same image. Therefore
\[
\gamma=u\alpha^{m-2}
\]
is the identity on all but the last core position, which is mapped to its predecessor.
\end{proof}

\begin{proposition}[Synchronization]\label{prop:sync}
Every automaton in the construction is synchronizing.
\end{proposition}
\begin{proof}
By Lemma \ref{lem:entry}, first map $Q$ into $C$. By Lemma \ref{lem:merge}, the semigroup restricted to $C$ contains an adjacent idempotent merge. Conjugating $\gamma$ by powers of the full cycle $\alpha$ moves this adjacent merge to every edge of the cycle. Applying these merges successively reduces $C$ to one state. Prepending the entry word therefore yields a reset word for $Q$.
\end{proof}

\section{The full-set boundary}
Suppose $S=Q$. A synchronizing automaton must contain a singular letter: if every letter were a permutation, every word would be a permutation and no reset word could exist. A singular letter compresses $Q$ in one step, while the empty word is not successful. Hence
\[
\tauac(\Acal;Q,A)=1.
\]
This proves the remaining part of Theorem \ref{thm:main}.

\section{Computational audit}
The proof is analytic, but we also implemented an independent verifier. For each pair
\[
1\le r\le30,\qquad 3\le L\le40,
\]
the program constructs the automaton, performs breadth-first search in the subset automaton to obtain the exact first-success distance, checks strong connectivity, and evaluates an explicit reset word built from the core entry and conjugated adjacent merges. This gives 1,140 tested parameter pairs, up to $n=70$. Every case passed:

\begin{center}
\begin{tabular}{lr}
\toprule
Check & Result\\
\midrule
Parameter pairs & 1,140\\
Exact distance $=L$ & 1,140 / 1,140\\
Strong connectivity & 1,140 / 1,140\\
Explicit reset word & 1,140 / 1,140\\
Failures & 0\\
\bottomrule
\end{tabular}
\end{center}

The verifier uses no external packages and is intended as reproducibility support rather than as a substitute for the proof.

\section{Discussion}
The second diagonal is the first infinite proper-subset line on which the avoid-or-compress upper bound can be tested against simultaneously binary, strong-connectivity, and synchronization constraints. Theorem \ref{thm:main} shows that none of these restrictions improves the general value: the full $n-|A|$ delay remains attainable.

This also isolates where more interesting behavior can first occur. On the third diagonal $|S|-|A|=3$, the classical subset-compression bound already beats $n-|A|$ when $|S|=n-1$, creating a boundary phase transition. That case is treated separately in the companion paper on the third diagonal.

\section*{Data and code availability}
The accompanying script \texttt{NO17\_SECOND\_DIAGONAL\_VERIFY.py} reconstructs the family and reproduces the computational audit reported above.

\begin{thebibliography}{9}
\bibitem{Szykula2018}
M. Szyku\l a,
\newblock Improving the Upper Bound on the Length of the Shortest Reset Word,
\newblock \emph{35th Symposium on Theoretical Aspects of Computer Science (STACS 2018)}, LIPIcs 96, Article 56, 2018.
\newblock doi:10.4230/LIPIcs.STACS.2018.56.

\bibitem{Szykula2026}
M. Szyku\l a,
\newblock Synchronizing Automata: Open Problems,
\newblock \emph{Electronic Proceedings in Theoretical Computer Science} 451 (2026), 33--47.
\newblock doi:10.4204/EPTCS.451.3; arXiv:2608.24245.

\bibitem{Volkov2022}
M. V. Volkov,
\newblock Synchronization of finite automata,
\newblock \emph{Russian Mathematical Surveys} 77(5) (2022), 819--891.

\bibitem{Ferens2021}
R. Ferens, M. Szyku\l a, and V. Vorel,
\newblock Lower Bounds on Avoiding Thresholds,
\newblock \emph{MFCS 2021}, LIPIcs 202, Article 46, 2021.
\newblock doi:10.4230/LIPIcs.MFCS.2021.46.
\end{thebibliography}

\end{document}
